\section{Továbbfejlesztési lehetőségek}

A WatchDB jelenlegi változata az fő MVP funkciókra koncentrál, ezért a tovább fejlesztési lehetőségek megtalálhatóak az további MVP verziókban. Az egyik legfontosabb továbbfejlesztési lehetőség a közösségi funkciók bevezetése. Ide tartozhatnak a publikus profilok, más felhasználók követése, valamint egy aktivitási feed, ahol láthatóvá válik, hogy a felhasználó ismerősei milyen filmeket néztek meg vagy értékeltek.

További fejlesztési irány lehet a személyes statisztikák és elemzések megjelenítése. A felhasználók számára hasznos lehet, ha áttekinthetik, milyen műfajú filmeket néznek leggyakrabban, hogyan változik az értékelési szokásuk, vagy egy adott időszakban mennyi filmet rögzítettek.

A filmek kezelését tovább bővíthetné az egyedi listák létrehozása. A jelenlegi watchlist és watched log mellett a felhasználók saját listákat is készíthetnének, például kedvenc filmekből, tematikus válogatásokból vagy későbbi filmestekhez kapcsolódó gyűjteményekből. Ezek a listák később megoszthatóvá is válhatnának.

A keresési élmény is fejleszthető lenne. A jelenlegi cím alapú keresés mellett hasznos lehetne műfaj, év, népszerűség vagy értékelés szerinti szűrés és rendezés. Ez különösen akkor válna fontossá, ha az alkalmazás már nemcsak konkrét filmek megtalálására, hanem új filmek felfedezésére is szolgálna.

Technikai oldalról a tesztelési rendszer is tovább bővíthető. A jelenlegi Vitest alapú tesztek mellé később Playwright alapú end-to-end tesztek is készülhetnének, amelyek teljes felhasználói folyamatokat ellenőriznek böngészőben. Emellett hasznos lehetne coverage riport bevezetése, valamint külön tesztadatbázis használata az adatbázishoz kötött folyamatok megbízhatóbb ellenőrzésére.

Hosszabb távon a projekt mobilos használhatósága és teljesítménye is fejleszthető. Bár az alkalmazás webes felületként készült, a filmkövetés gyakran mobil eszközön történik, ezért a reszponzív felület finomítása, a betöltési sebesség javítása és akár PWA jellegű funkciók bevezetése is indokolt lehet. Ezek a fejlesztések hozzájárulhatnának ahhoz, hogy a WatchDB ne csak működő MVP, hanem hosszabb távon is kényelmesen használható filmkövető platform legyen.
